\medskip\noindent
The construction of the underlying sheaf of abelian groups does not
require the coefficient rings to be commutative. We spell this out.

\begin{definition}
\label{definition-hom-left-modules}
Let $X$ be a topological space and let $\mathcal{A}$ be a sheaf of
not necessarily commutative rings with $1$ on $X$.
A {\it sheaf of left $\mathcal{A}$-modules} is a sheaf of abelian
groups $\mathcal{F}$ endowed with a morphism
$$
\mathcal{A} \times \mathcal{F} \longrightarrow \mathcal{F}
$$
such that $\mathcal{F}(U)$ is a unital left
$\mathcal{A}(U)$-module for every open $U \subset X$.
A morphism of sheaves of left $\mathcal{A}$-modules is a morphism
of the underlying sheaves of abelian groups compatible with these
actions.

Let $\mathcal{F}$ and $\mathcal{G}$ be sheaves of left
$\mathcal{A}$-modules. We denote the abelian group of morphisms by
$$
\Hom_{\mathcal{A}}(\mathcal{F}, \mathcal{G}).
$$
If $T \subset X$ is any subset, endowed with its subspace topology,
and $i_T : T \to X$ is the inclusion, we set
$$
\Hom_{\mathcal{A}}(T; \mathcal{F}, \mathcal{G})
=
\Hom_{i_T^{-1}\mathcal{A}}
(i_T^{-1}\mathcal{F}, i_T^{-1}\mathcal{G}).
$$
The inverse image sheaves have their natural left
$i_T^{-1}\mathcal{A}$-module structures. If $T$ is open, this agrees
with
$$
\Hom_{\mathcal{A}|_T}
(\mathcal{F}|_T, \mathcal{G}|_T).
$$
\end{definition}

\begin{lemma}
\label{lemma-sheaf-hom-left-modules}
In the situation of
Definition \ref{definition-hom-left-modules}, the rule
$$
U \longmapsto
\Hom_{\mathcal{A}|_U}
(\mathcal{F}|_U, \mathcal{G}|_U)
$$
on open subsets of $X$ is a sheaf of abelian groups. Denote it by
$\SheafHom_{\mathcal{A}}(\mathcal{F}, \mathcal{G})$. For every open
$U \subset X$ we have
$$
\Gamma(U, \SheafHom_{\mathcal{A}}(\mathcal{F}, \mathcal{G}))
=
\Hom_{\mathcal{A}}(U; \mathcal{F}, \mathcal{G}).
$$

For every open $U \subset X$, put
$$
\mathcal{Z}_{\mathcal{A}}(U)
=
\left\{
z \in \mathcal{A}(U) \mid
z|_V a = a z|_V
\text{ for every open }V \subset U
\text{ and every }a \in \mathcal{A}(V)
\right\}.
$$
This defines a sheaf of commutative rings
$\mathcal{Z}_{\mathcal{A}}$ on $X$, and
$\SheafHom_{\mathcal{A}}(\mathcal{F}, \mathcal{G})$
is canonically a sheaf of
$\mathcal{Z}_{\mathcal{A}}$-modules. In particular,
$$
\Hom_{\mathcal{A}}(\mathcal{F}, \mathcal{G})
$$
is canonically a module over
$\Gamma(X, \mathcal{Z}_{\mathcal{A}})$, and there is a canonical
inclusion
$$
\Gamma(X, \mathcal{Z}_{\mathcal{A}})
\subset
Z(\Gamma(X, \mathcal{A})).
$$
\end{lemma}

\begin{proof}
Addition of morphisms of left $\mathcal{A}|_U$-modules is defined
pointwise, so the displayed groups are abelian. Suppose that
$U = \bigcup U_i$ and that
$$
\varphi_i :
\mathcal{F}|_{U_i} \longrightarrow \mathcal{G}|_{U_i}
$$
are compatible morphisms of left
$\mathcal{A}|_{U_i}$-modules. For an open $V \subset U$ and
$s \in \mathcal{F}(V)$, the sections
$$
\varphi_{i, V \cap U_i}(s|_{V \cap U_i})
\in
\mathcal{G}(V \cap U_i)
$$
agree on overlaps. They glue to a unique section of
$\mathcal{G}(V)$. These glued sections are compatible with
restriction and define a morphism
$\varphi : \mathcal{F}|_U \to \mathcal{G}|_U$.
Additivity and $\mathcal{A}|_U$-linearity can be checked after
restriction to the covering $\{U_i\}$. The same local argument gives
uniqueness. This proves the sheaf assertion and the formula for
sections.

Restrictions preserve the defining condition for
$\mathcal{Z}_{\mathcal{A}}$. Conversely, suppose that
$U = \bigcup U_i$ and compatible sections
$z_i \in \mathcal{Z}_{\mathcal{A}}(U_i)$ glue to
$z \in \mathcal{A}(U)$. If $V \subset U$ is open and
$a \in \mathcal{A}(V)$, then
$$
z|_V a = a z|_V
$$
can be checked after restriction to the covering
$\{V \cap U_i\}$ of $V$. Thus $z$ belongs to
$\mathcal{Z}_{\mathcal{A}}(U)$. Hence
$\mathcal{Z}_{\mathcal{A}}$ is a sheaf of rings. Its sections are
closed under the ring operations, and two of its sections commute
by the defining condition. Thus it is a sheaf of commutative rings.

Let $z \in \mathcal{Z}_{\mathcal{A}}(U)$ and
$\varphi \in
\SheafHom_{\mathcal{A}}(\mathcal{F}, \mathcal{G})(U)$.
For $V \subset U$ open, define
$$
(z\varphi)_V(s) = z|_V\varphi_V(s).
$$
If $a \in \mathcal{A}(V)$ and $s \in \mathcal{F}(V)$, then
$$
(z\varphi)_V(as)
=
z|_V a\varphi_V(s)
=
a z|_V\varphi_V(s)
=
a(z\varphi)_V(s).
$$
Thus $z\varphi$ is a morphism of left
$\mathcal{A}|_U$-modules. The module axioms and compatibility with
restriction follow pointwise. Taking $U=X$ gives the asserted global
action. Finally, a locally central global section commutes with every
global section, which gives the displayed inclusion.
\end{proof}

\medskip\noindent
The last inclusion need not be an equality. For example, let
$X = \{x, u\}$ have open subsets
$\emptyset$, $\{u\}$, and $X$. Given a field $k$, the data
$$
\mathcal{A}(X) = k[t], \qquad
\mathcal{A}(\{u\}) = M_2(k), \qquad
t \longmapsto e_{11}
$$
define a sheaf of rings. The section $t$ lies in
$Z(\mathcal{A}(X))$, but its restriction $e_{11}$ is not central.
Thus $t$ does not belong to
$\Gamma(X, \mathcal{Z}_{\mathcal{A}})$.
